% !TEX root = ../main.tex
\chapter{Example workflows}\label{ch:workflows}
This chapter walks through complete pieces of research done with the app, one per section.
The first is the path the researcher and the supervisor share: from spikes to statistics
about how the spikes are spaced. The other six are the project's six research questions, in
the order of the research files that track them (\code{docs/rq_roundA/}). Each section says
what the question is, the clicks, what you should see, how to judge it, and what the app
cannot yet answer.

\begin{plainterms}[A workflow is a recipe, not a result]
A cookbook tells you how to bake the cake; it cannot promise yours will rise. Each section
below is the recipe. Where the project has tried it on real recordings, the section says
what came out and how sure anyone can be. Where nobody has, it says so.
\end{plainterms}

\begin{watchout}[The status of each question moves]
What each question can do is recorded, with dates, in its own file under
\code{docs/rq_roundA/}, and the clicks in \code{docs/RESEARCH_RUNBOOK.md}. This chapter
follows those files as they stood on 15~October~2026. Where they disagree with the app in
front of you, trust the app and the file, in that order.
\end{watchout}

%% ------------------------------------------------------------------------------------------
\section{From spikes to spike-train statistics}\label{sec:workflows-spikes}
The question: does a channel produce \term{spike}{spikes}, and if it does, are they spaced
at random, evenly, or in bursts? The answer comes in three parts: find the spikes, measure the
gaps between them, and ask whether the result beats chance. This is the path both the
researcher and the supervisor read, so it is the one to run most carefully.

\begin{plainterms}[Spikes, trains and gaps]
Think of raindrops on a tin roof. Each drop is a spike. The time between two drops is an
\term{inter-spike-interval}{inter-spike interval}. A shower that comes in gusts, several drops
quickly and then a pause, is made of \term{spike-train}{trains}. Whether the drops fall
steadily, at random or in gusts is what the interval statistics tell you.
\end{plainterms}

\subsection*{Step 1: pick a span}
\begin{enumerate}
\item Open \apppage{explore.signal} on the channel you want.
\item Drag across the trace to mark a \term{span}{span}. A few hours is a good start: long
  enough to hold many spikes, short enough to run in seconds.
\item Press \ui{Send span to Analyse}.
\end{enumerate}

\subsection*{Step 2: detect the spikes}
On \apppage{analyse.chain} you have three ways to find spikes. Choose one; they rest on
different ideas of what a spike is.
\begin{itemize}
\item \textbf{The simple rule.} \ui{Insert a stage} and pick \blockref{detection.spike_v1}. It
  marks a spike wherever the signal stands out from its own local average, up or down, by
  more than \emph{delta\_mv} millivolts; candidates closer together than \emph{d} samples
  count as one spike, at the most extreme of them. It starts from the rule of
  \citet{adamatzky2023multiscalar}, with changes its block entry lists.
\item \textbf{The Dehshibi detector.} \ui{Import} the \code{dehshibi_spikes} template. It is
  three blocks, a wavelet transform, a summation and a threshold, that follow the authors'
  own code for \citet{dehshibi2021electrical}. Each stage's output is drawn on its block page,
  so you can see what was kept and what was dropped. The older one-block version,
  \blockref{detection.dehshibi_spikes}, is kept only so old recipes still run.
\item \textbf{The inverted drop detector.} \ui{Import} the \code{spike_event_features}
  template. It turns the signal upside down, runs the project's drop detector, measures each
  event's shape (told that the signal was inverted, so the events are reported as spikes) and
  ends with \blockref{interrogation.intervals}.
\end{itemize}
Press \ui{Run chain}. Open the detector's block page from its row and look at the spikes over
the trace before you trust any number built on them.

\begin{watchout}[Three things that make spike counts wrong]
\begin{itemize}
\item \emph{Units and noise.} \blockref{detection.spike_v1} reads \emph{delta\_mv} in
  millivolts and converts it by the recording's declared unit. A recording with no declared
  unit is read as volts, and the chain row says ``read as volts''; if the file is really in
  millivolts, the threshold is then a thousand times too small and the block fires on
  noise. Check the recording's unit in \apppage{settings.datasets} first. The 1\,mV default
  is a placeholder: on a noisier channel raise it until the readout on the block page is no
  longer amber.
\item \emph{Samples, not seconds.} \emph{width} and \emph{d} count samples, so the same
  setting means ten times less time on a 10\,Hz recording than on a 1\,Hz one.
\item \emph{Bursts become one spike.} \blockref{detection.spike_v1} joins candidates that
  are each closer than \emph{d} samples to the last into one spike, however long the run.
  Spikes inside a tight burst are not counted separately, so the intervals inside it are
  never measured. Lower \emph{d} if your spikes come closer together than that.
\end{itemize}
\end{watchout}

\subsection*{Step 3: trains and intervals}
\begin{enumerate}
\item If the chain does not already end in it, \ui{Insert a stage} after the detector and pick
  \blockref{interrogation.intervals}. It measures the gap before and after every event and,
  for each group of events, the median gap, how variable the gaps are (CV), whether they drift
  over time (a trend and a drift ratio), and two local measures, LV and CV2, which compare each
  gap only with the next. Near 0 means clockwork, near 1 means random. The block prints the
  rule behind every number beside it.
\item To find bursts, insert \blockref{interrogation.trains} between the detector and
  Intervals. It sorts gaps into ``short, inside a train'' and ``long, between trains'' by
  looking for two humps in a histogram of the gaps. When there is one hump, it says ``no
  trains'' rather than inventing some. Turn on \emph{write\_group} if you want Intervals
  to measure timing \emph{within} each train.
\end{enumerate}

\begin{watchout}[Trains is experimental]
Trains found nothing beyond its surrogate null on the real recordings it was tried on
(\code{docs/prompts/blocks/reports/W1-evaluation.md}), and its template
\code{event_trains} is kept in the code but not offered in the app. ``No trains'' can also
mean that the gaps inside and between trains are too alike to tell apart. LV and CV2 are
very noisy on short runs: a group with fewer than three events gets none, with the reason.
\end{watchout}

\subsection*{Step 4: judge the count against chance}
\begin{enumerate}
\item On the Analyse toolbar, turn \ui{surrogate null} on. It is off by default, because
  tuning is meant to be fast.
\item Press \ui{Run chain} again. After the real run, the app runs the whole chain again on
  surrogate copies of the span: 20 by default, set in \apppage{settings.nulls}.
\item Read the footer. It says, for example, ``detected 120 · null expects 31 over 20 phase
  randomize draws · 3.9× null''. A count well above what the null expects is evidence that
  the detector is finding something the noise does not produce.
\end{enumerate}
To make the same comparison on many channels at once, save the detector alone as a template
(\ui{Save template}) and apply it in \apppage{discovery.runs}, which always draws a null for
every channel. Discovery's \ui{Apply template} lists only detection templates; a chain that
ends in Intervals or Trains is an interrogation template and is not offered there.

\subsection*{What this does not yet tell you}
The null above judges the \emph{number} of spikes. It does not judge the interval
statistics: no part of the app yet compares an LV, a CV or a train count with what shuffled
spike times would give. That null was specified and never built
(\code{docs/prompts/fixup/future/R-interrogation-null.md}). Until it is, read the interval
numbers as a description of the spikes you found, and say so when you report them. Nor does
any test in the repository check either spike detector against spikes a person marked by
hand. \Cref{app:nulls} explains the nulls in full.

%% ------------------------------------------------------------------------------------------
\section{Do cluster labels train a better classifier than hand labels?}\label{sec:workflows-rq1}
The question (RQ1): you can teach a computer to sort stretches of signal in two ways. You can
show it stretches you sorted yourself, \emph{interesting} or \emph{not interesting}. Or you can
let a clustering method group similar-looking stretches first and teach the computer those
groups. Give both the same exam on data neither has seen: which does better?

\begin{plainterms}[Two ways to learn birds]
One student learns birds from a field guide a person wrote. Another sorts a thousand
photographs into piles by look, names the piles, and learns from those. Then both sit the same
exam. The question is not which pile-maker is cleverer; it is whose knowledge carries over to
birds neither has seen.
\end{plainterms}

The question has two forms in the app. The first, the \emph{baseline}, is finished and has a
result. The second, which the researcher made the main one on 5~October~2026, is built and
waiting for the researcher's own blind labels.

\subsection*{The baseline: a paired job}
\begin{enumerate}
\item In \apppage{models.launch}, choose a training template, the recording and its channels,
  and the two label arms: A, the human labels, and B, one clustering of the pooled training
  \term{window}{windows}.
\item Under \emph{Evaluation} the split is blocked by time, with a gap of at least one window
  between training and test (the windows overlap, so a random split would leak; see
  \term{blocked-split}{blocked split}).
\item Press \ui{Train locally}. Both arms train the same random forest; a label-shuffle null
  is drawn for each.
\item Read \apppage{models.results} and \apppage{models.compare}: the difference in macro F1
  with its interval, McNemar's test, and results per class and per channel.
\end{enumerate}
\screenshot{pages/zz_models_ab/models.compare--fixup-ab-a-vs-b}{Models › Compare puts the two
arms side by side on the same exam windows: one difference, the label source, is changed.}

\emph{What came out} (run 79, real database, recorded in
\code{docs/rq_roundA/RQ1-cluster-vs-manual-labels.md}): on channels never trained on, the
hand-label arm scored macro F1 0.793 and the cluster arm 0.635; the difference, $+0.159$
with 95\% interval $[0.120, 0.197]$, excludes zero. The hand labels won on both exams, as the
researcher predicted before the run. The research file explains why this comparison is tilted
towards the hand labels: the clusters are translated into \emph{interesting} or not by
majority vote on the same hand labels, so the cluster arm can tie or lose but hardly win.

\subsection*{The main form: cluster the unlabelled pool, check blind}
\begin{enumerate}
\item In Library › Window sets, \ui{New window set} cuts a recording into non-overlapping
  windows at 1, 10 and 30 minutes, ignoring the hand labels.
\item In \apppage{analyse.chain}, \ui{Start from a Window pool}. The chain is Window pool
  $\to$ Trace shape $\to$ Shape clustering. A whole pack of channels is held out as an exam.
  The Shape clustering block's page is a dendrogram: drag the cut, look at each cluster's most
  typical member and a dozen others, and say which clusters are interesting.
\item \ui{Train model} opens Models › Launch filled in, with arm \emph{B.2 cluster labels ·
  trace shape}. \ui{Train locally} trains a random forest on the clusters. A CNN arm is
  written as a cluster script (\ui{Create SLURM script}) and brought back through
  \apppage{jobs.all}.
\item In \apppage{models.results}, \ui{Label test windows blind} opens a Review queue that
  shows each test window's raw trace with nothing else: no prediction, cluster, channel or
  earlier label. Answer \ui{I}, \ui{N}, \ui{U} (can't tell) or \ui{A}. Some windows come back
  twice, so your own consistency is measured too. Your first label freezes the cut and the
  mapping, so they cannot be tuned to the answers.
\item Read \emph{against a blind human} in Results: agreement per exam with its interval, the
  label-shuffle null, and the existing hand-label models scored on the same windows as a
  comparison line.
\end{enumerate}

\subsection*{How to judge it}
Read each exam separately, never pooled. A difference counts when its 95\% interval excludes
zero, the rule the researcher fixed before the baseline. Agreement with a blind person is the
yardstick for the main form; the forest reproducing its own clusters is not a score.

\subsection*{What is not yet answerable}
\begin{itemize}
\item The main result needs the researcher's own blind labelling of a real run; every number
  from that path so far is a placeholder from a test walk.
\item The third exam, the \term{held-out}{held-out} recording, is refused by every page until
  after the freeze, by design. It is run once.
\item The CNN arm has not yet been run on the cluster.
\end{itemize}

%% ------------------------------------------------------------------------------------------
\section{Can one accepted example find the rest?}\label{sec:workflows-rq2}
The question (RQ2): you point at one wiggle you have checked yourself and say ``find me more
like this''. The app searches the recordings for look-alikes, and you check each one by hand.
Are most of what it finds things you also accept, and does it find more than it would by
chance?

\begin{plainterms}[Searching by example]
It is like showing a shop assistant a photo of the shoe you want. They come back with a stack.
Some are your shoe, some only look like it in the photo. You learn how good the assistant is
by checking every box, and you learn whether they were trying by comparing with a stack
picked off the shelf at random.
\end{plainterms}

\subsection*{The clicks}
\begin{enumerate}
\item In \apppage{review.inspector}, on a detection worth searching for, press \ui{S}. It
  becomes a Library entry (an \term{exemplar}{exemplar}); the message names the entry.
  Alternatively, in \apppage{explore.signal}, mark a span and take it for Review: it appears as
  a seed under \emph{Explore selection}.
\item Press \ui{Seed search in Discovery}, or open \apppage{discovery.seed} and choose the
  seed there. The picker reaches every Library entry and filters by how it was made: your own
  first.
\item Set the \ui{Scope}: recording, channels and a section of hours. Keep it to hours, not the
  whole recording, while you learn.
\item Press \ui{Run seed search}. The page draws a histogram of how close each match is, with
  the matches the null gives behind it, and a recommended cut. Drag the cut and run again.
\item \ui{Open in Runs}, then \ui{Send \(n\) unjudged to Review}, and judge every match.
\item Back on Runs, \ui{Refresh after reviewing} and read the seed run's row.
\end{enumerate}
\screenshot{pages/discovery/discovery.seed--done}{The Seed page after a search: the real
matches and, behind them, what the null gives at each distance, with the cut drawn.}

\subsection*{What you should see, and how to judge it}
The row reads, in order, \emph{found}, \emph{already judged}, \emph{reviewed},
\emph{interesting}, \emph{precision}, \emph{null expects / draws} and \emph{× null}. The
answer is that row read with its null: precision says what share of the matches you accept,
and \emph{× null} says how many more matches the real signal gives than the surrogates do at
the same cut. Nothing on the row claims to beat chance unless \emph{× null} does. Precision
comes as two figures, each with its rule printed: one over the reviewed windows a match falls
in, one over the event-shaped human labels.

A seed search draws 200 surrogates per channel by default (\apppage{settings.nulls}). The
recommended cut is worked out per channel, each against its own null, at the significance
level and multiple-comparison correction set there; the sentence beside the cut says which.
A search estimated at more than ten minutes is offered as a cluster job instead
(\ui{Create SLURM script}); its result is brought back with \ui{Import HPC result}.

\subsection*{What is not yet answerable}
\begin{itemize}
\item \emph{Recall} on the row still uses an older rule that gives 0 on window labels by
  construction. Read precision and × null, not recall.
\item The first answer at scale, a seed search over a whole recording on the cluster, judged
  by the researcher, has not yet been made. The research file's walks in a sandbox used
  placeholder verdicts.
\end{itemize}

%% ------------------------------------------------------------------------------------------
\section{Does a shape keep its identity when it is stretched or squashed?}\label{sec:workflows-rq3}
The question (RQ3): is a slow, stretched-out version of a \term{motif}{motif} ``the same
motif'' as a fast, squashed one? The Library assumes yes: it groups motifs after stretching
them all to one length. If that is wrong, the Library's \term{family}{families} are wrong, so
``no'' is a real finding, not a failure.

\begin{plainterms}[The same tune, faster]
A song played at a faster tempo is still recognisably the same song. A different song played
at the same tempo is not. The question is whether the fungal shapes behave like songs, where
speed does not change identity, or like spoken words, where squashing ``cat'' can turn it
into something else.
\end{plainterms}

\subsection*{The clicks}
\begin{enumerate}
\item Open a family in \apppage{library.family}. Prefer a trough family: for sharkfins, how
  long the shape really is remains an open question, and the page says so.
\item Press \ui{Seed search in Discovery}. On the Seed page, turn on the \ui{scale bank}: the
  exemplar is searched at several lengths (by default 0.8, 1 and 1.25 times its own), each
  with its own null.
\item Run, send the matches to Review, and judge them.
\item On \apppage{discovery.runs}, add the accepted matches to the family. Each accepted pair
  gets one \term{edge}{edge} per way of measuring distance.
\item Back on the family's page, read the scale read-out: for each length, how many matches,
  how many judged and how many accepted, with the shape distance that ignores length beside a
  control that does not.
\end{enumerate}
\screenshot{pages/library_edges/library.family--fixup-v-trough-exemplar-edges-and-readout}{The
scale read-out on a trough family: one row per length, the length-blind distance beside its
control. The numbers in this picture come from a test walk with placeholder verdicts.}

\subsection*{How to judge it}
Identity holds under scale if acceptance stays high across the lengths while the length-blind
distance stays small and the control grows. Read each length against its own null: at the
same cut a shorter copy finds chance matches about twice as easily as the native length did
in the sandbox walk. Since 15~October a fourth distance, MPdist, is written on each edge too;
it asks whether two shapes share short pieces, and it is not itself a test of scale.

\subsection*{What is not yet answerable}
\begin{itemize}
\item The researcher's verdicts on a handful of trough exemplars' banked searches. Until then
  the read-outs show that the instrument works, not what the answer is.
\item Which grouping to read: the stored grouping g-05 and Ward re-run on today's events agree
  only partly (adjusted Rand 0.57), so g-05's families need re-checking first.
\item The sharkfin's own length, which decides whether any sharkfin row can be read
  (\code{docs/prompts/fixup/future/N-event-extent.md}).
\end{itemize}

%% ------------------------------------------------------------------------------------------
\section{Do bands and symbols find what the raw signal hides?}\label{sec:workflows-rq4}
The question (RQ4): split the recording into slow, middle and fast layers, turn each layer
into a string of letters (up, down, flat), and search the letters for patterns. Does that
find regions that a method working on the plain signal misses? You check the extra finds by
hand to see whether they are real.

\begin{plainterms}[Bass, middle and treble]
A recording is like a song with bass, middle and treble mixed together. A quiet bass line can
be lost under the treble. Turn the treble down and you may hear it. The band filter is the
mixing desk; the letters are a way to write down the tune so a computer can search it.
\end{plainterms}

\subsection*{The clicks}
\begin{enumerate}
\item In \apppage{discovery.runs}, set the \ui{Scope}: a recording, its channels and a section.
\item \ui{Apply template} with a raw-signal detector and no band: \code{drop_detection_v1},
  or the plain amplitude cut \code{signal_threshold}. This is side A.
\item \ui{Apply template} again with the symbolic template \code{symbol_search}, this time
  ticking bands under \emph{Bands}: the named bands from \apppage{settings.analysis-defaults}
  (0.001--0.01, 0.01--0.1 and 0.1--0.45\,Hz unless edited) and, if you like, \emph{Wavelet
  levels}. Each band becomes one run across the channels, with its own null.
\item Apply \code{symbol_search} once more with no band, for the like-for-like comparison.
\item Tick side A and the band set as side B, then \ui{Compare}. The overlap reads
  \emph{only A}, \emph{both} and \emph{only B}, split by verdict.
\item \ui{Send only-B unjudged to Review}, judge them, then \ui{Refresh after reviewing}.
\end{enumerate}
\screenshot{pages/band_scope/discovery.compare--fixup-z-band-set-as-side-b}{Compare with every
band run of one application taken together as side B, one row per band beneath.}

\subsection*{How to judge it}
The answer on that scope is the sentence under the verdict split: of the regions only a band
found, how many a person has judged and how many they accepted, with the bands' null count
beside it. Then use the like-for-like comparison to see how much is due to the band alone
rather than to the symbolic method.

\begin{watchout}[So far the null out-finds the signal]
On the 3-channel, 4-hour scope the project measured, the surrogates gave more regions than
the real signal on every band, wavelet layers included (for example 75 found against 247.75
expected on 0.01--0.1\,Hz, 20 draws per band). At those untuned settings nothing yet beats
chance. Counts stored before the symbolic encoder was given a fixed seed (3~October) should
be re-run once with \emph{force} before you read them.
\end{watchout}

\subsection*{What is not yet answerable}
Nothing the question needs is still to be built. Two gaps remain in running it: a band run
sent to the cluster comes back with no null, and the cluster-script writer writes a script
for the first pending run only, so a large band application gets one script, not one per
band. The answer itself needs the researcher's verdicts.

%% ------------------------------------------------------------------------------------------
\section{Where do people and the machine disagree?}\label{sec:workflows-rq5}
The question (RQ5): sometimes you say ``interesting'' where the detector saw nothing, and
sometimes the detector fires where you would say ``not interesting''. Where do those
disagreements happen? Do they bunch up on certain channels, times or shapes, or are they
scattered?

\begin{plainterms}[Two markers, one exam]
Two teachers mark the same essays. Where they disagree tells you something: if the
disagreements all fall on essays about one topic, one of them has a blind spot there. If they
are scattered, it is ordinary noise.
\end{plainterms}

\subsection*{The clicks}
\begin{enumerate}
\item In \apppage{discovery.runs}, \ui{Apply template} over a scope. Each run carries its
  null.
\item \ui{Send \(n\) unjudged to Review} and judge. Each verdict on a detection is an
  \term{adjudication}{adjudication}, which fills the ``machine yes, human no'' direction.
\item Open \apppage{discovery.compare} with \emph{human annotations} as one side. It shows
  four cells (machine yes or no, against human yes or no) and a fifth, \emph{not comparable},
  for stretches no run covered or no person reviewed; those are never counted as
  disagreements. Beneath, the two disagreement cells are broken down by channel, by time and
  by the tags people gave.
\item In \apppage{explore.corpus}, \emph{colour by · disagree} draws the same two cells over
  the map of the whole recording.
\end{enumerate}
\screenshot{pages/zzz_x_divergence/discovery.compare--fixup-x-human-side-pairs-by-the-divergence-with-the-breakdown}{Compare
against the human record: the four cells, \emph{not comparable} kept apart, and the
breakdown of the disagreements.}

\subsection*{How to judge it}
Look at whether the disagreement cells concentrate on a few channels, hours or shapes. The
page says plainly that it has no test of ``structured'' and no null for it: whether the
bunching is more than chance is, for now, your judgement, and should be reported as such.

\subsection*{What is not yet answerable}
\begin{itemize}
\item There is too little data. On 15~October the research file recorded no adjudications in
  the real database, and only a few dozen detections inside a reviewed window at all. The
  instrument is built; the hours in Review are not.
\item The machine side has no shape classes: no detector stores one, and classes given in
  Review are not yet stored.
\item Recall still uses an older rule that gives 0 on window labels by construction.
\end{itemize}

%% ------------------------------------------------------------------------------------------
\section{Does a pattern recur across channels and recordings?}\label{sec:workflows-rq6}
The question (RQ6): if the same motif shows up on many electrodes, that could mean three
things. One wire picking up another's signal (an \term{artifact}{artifact}): the copies sit at
exactly the same instant. The signal travelling through the mycelium (propagation): the copy
on the next electrode is a little later. Or the organism doing the same thing again in
different places or at different times. Once the first two are taken out, is there still
recurrence?

\begin{plainterms}[Echoes and choruses]
You hear a phrase sung from several rooms. If every room sings it at the very same moment,
it is probably one singer on a loudspeaker system. If each room sings it a moment after the
last, it is passing from room to room. If rooms sing it at different times of their own
accord, the phrase is something they all know. Only the last is recurrence.
\end{plainterms}

\subsection*{The clicks}
\begin{enumerate}
\item In \apppage{discovery.runs}, \ui{Apply template} across channels. Each channel gets
  its own surrogates, so a family built from the run has a null beside its counts.
\item Build the family's members and \term{edge}{edges} from the accepted matches (as in
  \cref{sec:workflows-rq3}).
\item On \apppage{library.family}, \ui{Classify across channels}. Each member is compared with
  every other channel of its recording over the same stretch of time. A pair counts only if it
  beats the same comparison at 100 random other times. The machine flags a \emph{suspected
  artifact} (copies within 1\,s at $|r| \geq 0.98$) and propagation (1--50\,s at
  $|r| \geq 0.5$); members under 30 samples are \emph{too short to tell}. These thresholds
  are settings in \apppage{settings.analysis-defaults}.
\item \ui{Send suspected artifacts to Review}. Only artifacts a person confirms are taken out.
\item Read \apppage{library.recurrence} with \emph{excluding artifacts} and
  \emph{propagation counted once}. A cell with no reviewed coverage reads \texttt{?}.
\end{enumerate}
\screenshot{pages/zzzz_ad_suspected_artifact/library.family--fixup-ad-classify-against-chance}{Classify
across channels on a family: each pair carries its chance test, and the machine only flags.}

\subsection*{How to judge it}
Recurrence persists if the counts with artifacts excluded and propagation counted once stay
above what the run's surrogates give. The Library shows only members above the dataset's
noise floor by default, so recurrence is counted over members large enough to be events.

\subsection*{What is not yet answerable}
\begin{itemize}
\item The three families tried first (from a 10\,Hz recording) cannot answer it: most of
  their members are too short and too small to tell. A family built from a run, on larger
  and longer members, can.
\item The 64 members flagged on M2\_aug wait for a person to judge them.
\item Families imported rather than found by a run have no surrogates, so their counts have
  no null; the page says so.
\item ``Across recordings'' of other species is only starting: recordings of other fungi
  were derived in October 2026 and some were registered by 15~October. Each needs its own
  noise floor set in \apppage{settings.datasets} before its members are counted.
\end{itemize}
